\documentclass[11pt,a4paper]{article}
\usepackage[margin=22mm,headheight=15pt]{geometry}
\usepackage[T1]{fontenc}
\usepackage{lmodern,amsmath,amssymb,graphicx,booktabs,tabularx,array}
\usepackage{xcolor,fancyhdr,listings,caption,hyperref}
\definecolor{ink}{HTML}{17364A}
\definecolor{accent}{HTML}{176B87}
\definecolor{light}{HTML}{F2F5F7}
\hypersetup{colorlinks=true,linkcolor=accent,urlcolor=accent,
  pdftitle={Modelling Annual Average Daily Traffic},pdfauthor={}}
\pagestyle{fancy}
\fancyhf{}
\fancyhead[L]{\small\textcolor{ink}{MH3510 \quad Regression Analysis}}
\fancyhead[R]{\small AADT project}
\fancyfoot[C]{\small\thepage}
\renewcommand{\headrulewidth}{0.3pt}
\setlength{\parindent}{0pt}
\setlength{\parskip}{6pt}
\setlength{\emergencystretch}{2em}
\renewcommand{\arraystretch}{1.2}
\captionsetup{font=small,labelfont=bf}
\newcolumntype{Y}{>{\raggedright\arraybackslash}X}
\newcommand{\fig}[3]{%
  \begin{center}\begin{minipage}{\linewidth}\centering
  \includegraphics[width=#2\linewidth]{assets/#1}
  \captionof{figure}{#3}\end{minipage}\end{center}}
\lstdefinestyle{rcode}{language=R,basicstyle=\ttfamily\footnotesize,
  keywordstyle=\color{accent},commentstyle=\color{gray},
  stringstyle=\color{ink},breaklines=true,breakatwhitespace=false,
  columns=fullflexible,keepspaces=true,showstringspaces=false,
  frame=single,rulecolor=\color{gray!30},backgroundcolor=\color{light},
  numbers=left,numberstyle=\tiny\color{gray},numbersep=7pt,
  xleftmargin=12pt,framexleftmargin=4pt,tabsize=2}
\lstdefinestyle{console}{basicstyle=\ttfamily\scriptsize,
  breaklines=true,breakatwhitespace=false,columns=fullflexible,
  keepspaces=true,showstringspaces=false,frame=none}
\newcommand{\AuthorName}{Eugene Ng}
\newcommand{\MatricNumber}{N22602219H}

\begin{document}
\hypersetup{pageanchor=false}
\begin{titlepage}
\vspace*{20mm}
{\large\color{accent} MH3510 \quad REGRESSION ANALYSIS\par}
\vspace{10mm}
{\Huge\bfseries\color{ink} Modelling Annual\\[3mm] Average Daily Traffic\par}
\vspace{6mm}
{\Large Multiple linear regression with R\par}
\vspace{18mm}
\begin{tabular}{@{}ll}
Name: & \AuthorName\\[6mm]
Matriculation number: & \MatricNumber\\[6mm]
Submission: & Individual project\\[3mm]
Due date: & 11 November
\end{tabular}
\vfill
{\large\bfseries Overview\par}
This study models traffic volume for 121 road sections using county population,
lane count, road width and access control. The initial model is checked using
residual plots and revised through transformations. All 11 candidate models,
including unsuccessful attempts, are retained.

A fifth-root response model with log population, log lane count and access control
provides a defensible final specification. For the requested road, its
back-transformed typical prediction is approximately \textbf{5,494 vehicles/day},
with an approximate 95\% individual prediction interval of
\textbf{1,164 to 17,928 vehicles/day}.
\vspace{8mm}
{\small The main report follows the five-stage modelling workflow.
Appendices contain diagnostic evidence, the complete R script and its output.}
\end{titlepage}
\hypersetup{pageanchor=true}

\section{Step 1: Data and graphical exploration}
Average annual daily traffic (AADT) is the average number of vehicles passing
through a road section each day over a year. The analysis uses only the first
five columns of \texttt{aadt.txt}; the steam dataset is not used.
\begin{center}
\begin{tabularx}{\linewidth}{@{}clY@{}}
\toprule Column & Symbol & Definition\\ \midrule
1 & $Y$ & AADT, in vehicles/day (response)\\
2 & $X_1$ & Population of the county containing the road section\\
3 & $X_2$ & Number of lanes\\
4 & $X_3$ & Road width, in feet\\
5 & $A$ & Access indicator: 1 for access control; 0 for no access control\\
\bottomrule
\end{tabularx}
\end{center}
The original fifth column uses 1 for access control and 2 for no control.
Set $A=1$ for original code 1 and $A=0$ otherwise. All 121 observations are complete:
27 have access control and 94 do not. AADT ranges from 201 to 155,547.

\fig{01_scatterplot_matrix.png}{0.69}{Scatterplot matrix of the response and four candidate predictors.}

Traffic has positive marginal associations with population, lanes and access
control, while width has a weaker association. Population and AADT are strongly
skewed. Lanes and access control are related: all two-lane sections lack access
control, while sections with more than four lanes have access control. The
regression estimates conditional associations, not causal effects.

\clearpage
\section{Step 2: Initial multiple linear regression}
The starting model includes all four required predictors:
\[
Y_i=\beta_0+\beta_1X_{1i}+\beta_2X_{2i}+\beta_3X_{3i}+\beta_4 A_i+\varepsilon_i.
\]
Ordinary least squares gives
\[
\widehat{Y}=-26035.84+0.033032X_1+9157.94X_2+100.29X_3+23610.36A.
\]
The initial model has $R^2=0.7527$, adjusted $R^2=0.7442$ and residual standard
error $15{,}292.6$. Its overall $F$-test is significant
($F(4,116)=88.29$, $p<2.2\times10^{-16}$).

\subsection*{Why this fit is insufficient}
The fitted relationship explains some variation, but statistical significance
does not establish adequacy. The residual-versus-fitted plot shows curvature and
increasing spread; the Q--Q plot has substantial tail departures. Moreover,
45 fitted AADT values are negative, although traffic counts cannot be negative.
These features motivate transformations rather than accepting the initial fit.

\fig{M1_raw_residuals.png}{1}{Diagnostics for M1, the untransformed four-predictor model.
Residuals are plotted against fitted values and every predictor; the Q--Q plot
checks approximate normality.}

The remedies follow Chapter 2, Part 4: transform predictors to address curvature,
transform the response to address changing spread, and refit and check again
\cite{chapter2}. The fitted coefficients and full console output are retained
in Appendix~\ref{app:output}.

\clearpage
\section{Step 3: Transformations and the model record}
The candidate sequence is exploratory. Logarithms compress high values strongly;
the square-root transformation compresses them less. Fourth and fifth roots are
intermediate positive-power trials. A fifth root is simply $Y^{1/5}$; fitting it
still uses ordinary \texttt{lm()} regression, linear in its coefficients.

The predictor sets used below are
$P_0=(X_1,X_2,X_3,A)$, $P_1=(\log X_1,X_2,X_3,A)$,
and $P_2=(\log X_1,\log X_2,X_3,A)$.
A reduced set omits width $X_3$. All logarithms are natural.

\begin{center}\small
\begin{tabularx}{\linewidth}{@{}llY@{}}
\toprule Model & Response; predictors & What improved or remained problematic\\ \midrule
M1 & $Y$; $P_0$ & Baseline: curvature, increasing spread, tail problems and negative fitted values.\\
M2 & $\log Y$; $P_0$ & Q--Q plot improves, but the residual trend curves downward at large fitted values.\\
M3 & $\log Y$; $P_1$ & Population transformation helps; curvature, uneven spread and tail observations remain.\\
M4 & $\log Y$; $P_2$ & Flatter mean trend, but residual spread decreases as fitted values increase.\\
M5 & $\log Y$; $P_1$ plus $X_2^2$ & Quadratic lanes improve M3, but do not resolve the changing spread.\\
M6 & $\log Y$; reduced $P_2$ & Width can be removed; the variance issue persists.\\
M7 & $Y^{1/4}$; $P_2$ & Much more stable spread and straighter Q--Q plot. A reasonable candidate.\\
M8 & $\sqrt{Y}$; $P_2$ & Some improvement, but curvature and increasing spread remain.\\
M9 & $Y^{1/4}$; reduced $P_2$ & Reasonable simpler fourth-root alternative; width is not supported.\\
M10 & $Y^{1/5}$; $P_2$ & Relatively flat residual trend, more stable spread and reasonable Q--Q plot.\\
M11 & $Y^{1/5}$; reduced $P_2$ & Selected working model: M10 simplified without losing a supported predictor.\\
\bottomrule
\end{tabularx}
\captionof{table}{All candidate models are retained, including unsuccessful fits.}
\end{center}

\subsection*{How the response power was chosen}
The main diagnostic justification is that both root models correct the pronounced
defects of the raw, log and square-root fits. M11 is a defensible working choice,
not a uniquely proven optimum; M9 remains a reasonable alternative.

For transparency, earlier exploratory Box--Cox profiling suggested a power near
0.18 and led to the nearby fifth-root trial. That extended work is archived.
It is not required in the course-level script or used here as a substitute for
the residual checks. The exact fifth root is an extension of the lecture's
transformation principle, permitted by the assignment, rather than a transformation
prescribed in the slides.

\textbf{Comparison rule.} $R^2$, adjusted $R^2$ and residual standard deviations
refer to the fitted response scale. They should not be ranked across different
response transformations as if the response were unchanged.
An $F$-test also cannot compare different transformed responses.

\clearpage
\section{Step 4: Simplification and the final fitted model}
For nested models with the same response, the extra-sum-of-squares test is
\[
F=\frac{(\operatorname{SSE}_{R}-\operatorname{SSE}_{F})/(p_F-p_R)}
        {\operatorname{SSE}_{F}/(n-p_F-1)}.
\]
Here $p_F$ and $p_R$ count predictors excluding the intercept. This is the
full/reduced-model procedure in Chapter 4 Part II \cite{chapter4ii}.

\begin{center}
\begin{tabularx}{\linewidth}{@{}Yrr@{}}
\toprule Nested comparison & $F$ & $p$-value\\ \midrule
Add width to reduced log model: M6 vs M4 & 0.395 & 0.531\\
Add quadratic lanes: M3 vs M5 & 9.165 & 0.00305\\
Add width to fourth-root model: M9 vs M7 & 0.723 & 0.397\\
Add width to fifth-root model: M11 vs M10 & 0.632 & 0.428\\
Add squared log predictors to M11 jointly & 0.736 & 0.481\\
\bottomrule
\end{tabularx}
\captionof{table}{Model comparisons. Preliminary log-model tests are interpreted
alongside their unsatisfactory residual diagnostics.}
\end{center}

At the 5\% level, there is insufficient evidence that width adds explanatory value
to the fifth-root model after accounting for population, lanes and access control.
The residual plot against the omitted width also shows no obvious remaining trend.
The additional curvature test does not support the two extra squared terms.

Let $Z=Y^{1/5}$. The selected model is
\[
Z_i=\beta_0+\beta_1\log X_{1i}+\beta_2\log X_{2i}+\beta_4 A_i+\varepsilon_i,
\]
with fitted conditional mean
\[
\boxed{\widehat{Z}=-2.739015+0.535189\log X_1
          +2.317245\log X_2+1.106497A.}
\]

\begin{center}
\begin{tabular}{@{}lrrrr@{}}
\toprule Term & Estimate & Std.\ error & $t$ & $p$-value\\ \midrule
Intercept & $-2.739015$ & 0.554451 & $-4.940$ & $2.63\times10^{-6}$\\
$\log X_1$ & 0.535189 & 0.052970 & 10.104 & $1.24\times10^{-17}$\\
$\log X_2$ & 2.317245 & 0.246455 & 9.402 & $5.58\times10^{-16}$\\
Access $A$ & 1.106497 & 0.213759 & 5.176 & $9.51\times10^{-7}$\\
\bottomrule
\end{tabular}
\captionof{table}{Final model coefficient estimates and individual $t$-tests.}
\end{center}

On the fifth-root scale, $R^2=0.8528$, adjusted $R^2=0.8490$ and
$s=0.7478$ with 117 residual degrees of freedom.
The overall test gives $F(3,117)=225.95$, $p<2.2\times10^{-16}$.
All retained predictors have evidence of a nonzero conditional association.

\clearpage
\subsection*{Final residual checks and interpretation}
\fig{M11_fifth_root_reduced_residuals.png}{1}{Final-model diagnostics. Residual spread
is substantially improved, with a relatively flat fitted-value trend. Tail
observations remain and must be discussed.}

The residual cloud is substantially more satisfactory than M1's, and the Q--Q
plot is reasonably straight apart from a few tail observations. The plots support
an approximate working model; they do not prove the assumptions.

Holding the other fitted predictors fixed, doubling population increases
fitted $Z$ by $0.535189\log 2\approx0.371$; doubling lane count increases it by
$2.317245\log 2\approx1.606$. Access control is associated with a 1.1065 higher
fitted $Z$. These are fifth-root-scale changes, not fixed changes in vehicles/day
or percentages of AADT. The associations do not establish causal effects.

\textbf{Influence.} Row 93 has the largest Cook's distance: AADT 756, population
941,411, two lanes, width 24 feet and no access control. Row 92 also has a large
positive residual. With no evidence of recording errors, all observations remain.
Omitting row 93 only for sensitivity changes the requested typical prediction
from 5,494 to 5,430, about 1.2\%; retained coefficient signs remain positive.

\textbf{Independence.} The file-order Durbin--Watson statistic is 1.7726.
However, no time or spatial ordering is documented, so it does not justify a
serial-correlation conclusion. Independence remains an assumption; possible
county clustering cannot be resolved from the supplied variable description.
The influence and row-order plots are retained in Appendix~\ref{app:diagnostics}.

\clearpage
\section{Step 5: Prediction for the requested road}
The requested setting is population $50{,}000$, three lanes, width 60 feet and
original control code 2. Since code 2 means no access control, use $A=0$.
Width is recorded but does not enter the final reduced model.

\begin{lstlisting}[style=rcode,numbers=none]
new_road <- data.frame(population = 50000, lanes = 3,
                       width = 60, access = 0)
predict(final_model, new_road, interval = "confidence")
predict(final_model, new_road, interval = "prediction")
\end{lstlisting}

With $\mathbf{x}_0=(1,\log 50000,\log 3,0)^\mathsf{T}$ and design matrix $X$,
the lecture's intervals on the fitted $Z$ scale are
\begin{align*}
\text{Mean-response CI:}\quad&
\widehat Z_0\ \pm\ t_{0.975,117}s
 \sqrt{\mathbf{x}_0^\mathsf{T}(X^\mathsf{T}X)^{-1}\mathbf{x}_0},\\
\text{Individual PI:}\quad&
\widehat Z_0\ \pm\ t_{0.975,117}s
 \sqrt{1+\mathbf{x}_0^\mathsf{T}(X^\mathsf{T}X)^{-1}\mathbf{x}_0}.
\end{align*}
The extra 1 accounts for the variability of an individual road section.

\begin{center}
\begin{tabularx}{\linewidth}{@{}Yrrr@{}}
\toprule Quantity & Fitted value & Lower 95\% & Upper 95\%\\ \midrule
Mean fifth-root AADT: confidence interval & 5.59736 & 5.40290 & 5.79183\\
Individual fifth-root AADT: prediction interval & 5.59736 & 4.10375 & 7.09098\\
Typical AADT and individual prediction interval & 5494.36 & 1163.86 & 17927.96\\
\bottomrule
\end{tabularx}
\captionof{table}{Prediction results. The last row is in vehicles/day; the first
two rows are in $(\text{vehicles/day})^{1/5}$.}
\end{center}

The final row raises the fitted value and individual interval endpoints to the
fifth power. Because this transformation is increasing, the individual interval
can be back-transformed in this way.

\textbf{Mean versus typical traffic.} The value $5.59736^5\approx5494$ is a
back-transformed typical prediction (approximately a conditional median under
the normal-error working model), not the arithmetic mean AADT.
Similarly, raising the mean-$Z$ confidence limits to the fifth power does not
give a confidence interval for arithmetic mean AADT. The mean confidence
interval is reported on its fitted scale, without claiming an original-scale
mean interval or introducing advanced retransformation methods.

No three-lane sections are observed, so the requested lane count interpolates
between two and four lanes. The wide individual prediction interval reflects
substantial road-to-road variation and relies on the model assumptions.

\section{Conclusion and limitations}
A transformed multiple regression addresses the strongest deficiencies of the
initial model. M11 retains population, lanes and access control; width is omitted
using the full/reduced $F$-test and residual checks. M9 is a reasonable alternative,
giving a typical prediction near 5,728. The analysis is exploratory and uses the
same dataset to select and assess the model. The reported textbook intervals
are conditional on the selected specification and approximate error assumptions.

\clearpage
\section*{Reproducibility and references}
The analysis uses base R and all 121 observations. From the project directory run:
\begin{lstlisting}[style=rcode,numbers=none]
Rscript analysis.R
\end{lstlisting}
This generates the complete output, model objects and plots in
\texttt{coursework\_results/}. The supplied LaTeX source package includes the
script, dataset, console output and plot assets. To compile the report from
that package's directory:
\begin{lstlisting}[style=console]
tectonic aadt_report.tex
\end{lstlisting}

Appendix~\ref{app:diagnostics} preserves the remaining diagnostic figures.
Appendix~\ref{app:code} contains the full executable R code, and
Appendix~\ref{app:output} reproduces the console output. Formula environment
memory addresses in that output are R implementation details, not statistical
results. Confidence regions and Bonferroni limits are omitted as instructed.

The archived extended exploration is not necessary to reproduce the main
course-level analysis. No observations or unsuccessful candidate models have
been removed to improve the reported fit.

\begin{thebibliography}{9}
\bibitem{tutorial} MH3510 course material.
\emph{Modeling linear regression model with R}, Sections 2--7.
Practical workflow, model checking, reduced-model tests and prediction.
\bibitem{chapter2} MH3510 course material.
\emph{Chapter 2: Simple linear regression}, Part 4, printed slides 48--59 and 64.
Local file: \texttt{MH3510\_Chapter2 (1).pdf}.
\bibitem{chapter4i} MH3510 course material.
\emph{Chapter 4: Multiple linear regression, testing and prediction}, Part I.
Local file: \texttt{MH3510\_Chapter4\_Part1.pdf}.
\bibitem{chapter4ii} MH3510 course material.
\emph{Chapter four: Multiple linear regression (Part II)}, PDF pages 12--31
and 35--36. Local file: \texttt{MH3510\_CH4part2(ppt).pdf}.
\bibitem{data} Course project dataset, \texttt{aadt.txt}, and assignment instructions.
Only the first five columns are used.
\end{thebibliography}

\clearpage
\appendix
\section{Remaining candidate diagnostics}\label{app:diagnostics}
M1 and M11 appear in the main report. The following figures preserve all other
candidate fits. Model specifications and decisions are listed in Section 3.
\fig{M2_log_y_residuals.png}{0.85}{M2: log response with untransformed predictors.
Improved Q--Q plot, but a curved residual trend remains.}
\fig{M3_log_y_log_population_residuals.png}{0.85}{M3: log response and log population.
Curvature and changing spread still motivate revision.}
\clearpage
\fig{M4_log_y_log_population_lanes_residuals.png}{0.98}{M4: log response, log population
and log lanes. The spread still differs across fitted levels.}
\fig{M5_log_y_quadratic_lanes_residuals.png}{0.98}{M5: a quadratic lane term improves
the log model's mean specification but does not fully resolve uneven spread.}
\clearpage
\fig{M6_log_y_reduced_residuals.png}{0.98}{M6: dropping width simplifies the log model
but does not repair its variance issue.}
\fig{M7_fourth_root_full_residuals.png}{0.98}{M7: the full fourth-root model is a
reasonable transformed candidate with substantially improved diagnostics.}
\clearpage
\fig{M8_square_root_full_residuals.png}{0.98}{M8: the square-root response model
retains curvature and increasing residual spread.}
\fig{M9_fourth_root_reduced_residuals.png}{0.98}{M9: the reduced fourth-root model
remains a reasonable alternative to the final specification.}
\clearpage
\fig{M10_fifth_root_full_residuals.png}{1}{M10: the full fifth-root model.
The reduced model M11 omits width after the nested comparison.}
\fig{02_final_influence_and_order.png}{1}{Final-model Cook's distance and residuals
in file order. File order is not established to represent time.}
\clearpage
\fig{03_final_transformed_scatterplots.png}{0.9}{Transformed response and predictors.
Width remains displayed to assess the omitted variable.}

\clearpage
\section{Complete R code}\label{app:code}
The code below is the executable \texttt{analysis.R}. Helper functions only format
output and save figures. Model fitting and inference use base R.
\lstinputlisting[style=rcode]{analysis.R}

\clearpage
\section{Complete R console output}\label{app:output}
The following is the complete output from the course-level analysis. The model
summaries, nested tests, sensitivity results and predictions are included.
\lstinputlisting[style=console]{R_output.txt}
\end{document}
